\chapter{Обменный Real-символ степени семь и исправление соглашения}

Предыдущий гейт правильно построил неограниченный класс расширения
Тёплица, но неверно перевёл его номер $KKO_1$ в ковариантную нумерацию
$KO_n$. Настоящий гейт сначала исправляет эту ошибку, а затем строит
точный входной класс для уже доказанного выхода $15\in KO_6$.

\section{Почему степень пять была ошибкой}

В обозначениях Бурна, Келлендонка и Ренни extension-класс представлен
неограниченным модулем с $Cl_{0,1}$ и записан как элемент $KKO_1$;
см. предложение 3.3 в \path{arXiv:1604.02337v3}. Однако в ковариантной
таблице вещественной $K$-теории, использованной в Bott-сравнении,
граничная карта короткой точной последовательности имеет вид
\begin{equation}
 \partial_n:KO_n(A/J)\longrightarrow KO_{n-1}(J).
 \label{eq:v5-ko7-boundary-convention}
\end{equation}
Явные формулы всех восьми карт в этой нумерации даны Бурсема и Лорингом
в \path{arXiv:1504.03284v3}. В частности, они отдельно подчёркивают карту
$\partial_5:KO_5\to KO_4$, что исключает прежнее сложение $5+1=6$.

Следовательно, вход для выхода в $KO_6$ должен иметь степень семь:
\begin{equation}
 \boxed{7-1=6\pmod 8.}
 \label{eq:v5-ko7-correct-degree}
\end{equation}

\begin{nogocriterion}[Запрет смешения двух нумераций]
Номер $KKO_1$ неограниченного extension-модуля нельзя без перевода
складывать с индексом ковариантной таблицы $KO_n$. Все дальнейшие
boundary-вычисления выполняются по правилу
$KO_n\to KO_{n-1}$.
\end{nogocriterion}

\section{Коэффициентный класс степени шесть}

Пусть
\begin{equation}
 B=M_{105}(\mathbb C)_{\mathbb R}.
\end{equation}
По вещественной инвариантности Мориты и таблице для комплексных чисел,
рассматриваемых как вещественная алгебра,
\begin{equation}
 KO_{2k}(B)\simeq\mathbb Z,
 \qquad
 KO_{2k+1}(B)=0.
 \label{eq:v5-ko7-coefficient-table}
\end{equation}
Предыдущий Bott-гейт доказал обменный класс
\begin{equation}
 \kappa_{15}:=[T_{\mathbb R}]=15\in KO_6(B),
 \qquad
 c_6(\kappa_{15})=(-15,+15).
 \label{eq:v5-ko7-coefficient-fifteen}
\end{equation}

\section{Канонический winding-класс}

Обозначим через $C^*_{\mathbb R}(\mathbb Z)$ вещественную групповую
$C^*$-алгебру целых чисел и через $u_{\mathbb R}$ её канонический
унитарный генератор. Его winding-класс имеет степень один:
\begin{equation}
 [u_{\mathbb R}]\in KO_1(C^*_{\mathbb R}(\mathbb Z)).
 \label{eq:v5-ko7-real-winding}
\end{equation}
Это не новый физический носитель: тот же генератор представлен
двусторонним сдвигом $Ue_n=e_{n+1}$, уже участвующим в неограниченном
цикле оператора числа.

Здесь существенно, что $C^*_{\mathbb R}(\mathbb Z)$ соответствует
окружности с инволюцией $z\mapsto\bar z$, а не алгебре вещественнозначных
функций на окружности с тривиальной инволюцией. В обозначениях примера
10.8 Бурсема--Шокета ядро оценки в фиксированной точке является
десуспензией $S^{-1}\mathbb R$; именно поэтому winding-компонента в
$KO_1$ целочисленна. Более точно,
\begin{equation}
 KO_1(C^*_{\mathbb R}(\mathbb Z))\simeq\mathbb Z\oplus\mathbb Z_2,
\end{equation}
и $[u_{\mathbb R}]$ выбирается из редуцированного свободного слагаемого
$\mathbb Z$, а не из точечного торсионного слагаемого.

Искомый символ определяется внешним произведением
\begin{equation}
 \xi_{15}
 =[u_{\mathbb R}]\boxtimes\kappa_{15}
 \in KO_7(C^*_{\mathbb R}(\mathbb Z)\otimes B).
 \label{eq:v5-ko7-symbol}
\end{equation}
Его степень равна $1+6=7$.

Существование целочисленной редуцированной компоненты можно проверить
независимо. Расщеплённая оценка в единице и формула десуспензии дают
\begin{equation}
 \widetilde{KO}_7(C^*_{\mathbb R}(\mathbb Z)\otimes B)
 \simeq KO_6(B)\simeq\mathbb Z.
 \label{eq:v5-ko7-reduced-group}
\end{equation}
Именно в этой компоненте лежит $\xi_{15}$.

\section{Произведение с расширением Тёплица}

Граничная карта является модульной относительно внешнего произведения.
Для базового winding-класса расширение Тёплица даёт
\begin{equation}
 \partial_1([u_{\mathbb R}])=\varepsilon[1],
 \qquad \varepsilon\in\{+1,-1\},
 \label{eq:v5-ko7-base-index}
\end{equation}
где знак определяется общей ориентацией сдвига. Поэтому
\begin{align}
 \partial_7(\xi_{15})
 &=\partial_1([u_{\mathbb R}])\boxtimes\kappa_{15}\\
 &=\varepsilon\kappa_{15}
 =\varepsilon15\in KO_6(B).
 \label{eq:v5-ko7-boundary-fifteen}
\end{align}
После выбора ориентации, согласованной с $PUP=S$, комплексификация
\eqref{eq:v5-ko7-boundary-fifteen} равна ранее вычисленной паре
$(-15,+15)$. При одновременном обращении winding-ориентации меняются оба
знака, но абсолютный класс и вес не меняются.

\begin{theorem}[Классовая факторизация обменного индекса]
Доказанный класс $15\in KO_6(B)$ является граничным образом
канонического символа степени семь
$\xi_{15}=[u_{\mathbb R}]\boxtimes\kappa_{15}$.
Следовательно, неограниченный Toeplitz-extension с оператором числа $N$
и коэффициентный KO6-класс образуют корректную Kasparov-факторизацию
класса пятнадцать с точностью до общей ориентации.
\end{theorem}

Нормированный коэффициент остаётся
\begin{equation}
 \frac{|\partial_7(\xi_{15})|}{105}=\frac{15}{105}=\frac17.
\end{equation}

\section{Что ещё не закрыто}

Классовая факторизация теперь закрыта, но отдельная единичная матричная
формула для представителя $\xi_{15}$ в унитарной модели $KO_7$ ещё не
записана. Она нужна как независимый контроль симметрий, но уже не может
изменить группу, boundary-образ или число пятнадцать. Также из
K-теоретического произведения не следуют физическое действие, конечная
энергия, масса или масштаб.

\begin{protocol}
\begin{enumerate}
 \item прежняя степень пять: \textbf{отозвана};
 \item ковариантная boundary-конвенция: \textbf{$KO_n\to KO_{n-1}$};
 \item правильная входная степень: \textbf{семь};
 \item обменный коэффициентный класс $\kappa_{15}$:
 \textbf{$15\in KO_6(B)$};
 \item winding-класс: \textbf{$[u_{\mathbb R}]\in KO_1$};
 \item символ $\xi_{15}$: \textbf{построен как внешнее произведение};
 \item группа редуцированного символа: \textbf{$\mathbb Z$};
 \item boundary-образ: \textbf{$\pm15$};
 \item комплексификация: \textbf{$(-15,+15)$ с точностью до ориентации};
 \item нормированный вес: \textbf{$1/7$};
 \item классовая неограниченная факторизация: \textbf{закрыта};
 \item единичная матричная KO7-модель: \textbf{не записана};
 \item физическое действие: \textbf{не получено};
 \item следующий гейт: \textbf{явный унитарный представитель KO7}.
\end{enumerate}
\end{protocol}

Машинный сертификат сохранён в
\path{s2t_v5_real_toeplitz_degree_seven_symbol_gate_results.json}.